
\documentclass{article}
\usepackage{colortbl}
\usepackage{makecell}
\usepackage{multirow}
\usepackage{supertabular}

\begin{document}

\newcounter{utterance}

\centering \large Interaction Transcript for game `cladder', experiment `full\_v1.5\_default', episode 188 with reflexion{-}Qwen3.5{-}9B.
\vspace{24pt}

{ \footnotesize  \setcounter{utterance}{1}
\setlength{\tabcolsep}{0pt}
\begin{supertabular}{c@{$\;$}|p{.15\linewidth}@{}p{.15\linewidth}p{.15\linewidth}p{.15\linewidth}p{.15\linewidth}p{.15\linewidth}}
    \# & \multicolumn{2}{c}{Player} && \multicolumn{2}{c}{Game Master} \\
    \hline

    \theutterance \stepcounter{utterance}  
    & & & \multicolumn{4}{p{0.6\linewidth}}{
        \cellcolor[rgb]{0.9,0.9,0.9}{
            \makecell[{{p{\linewidth}}}]{
                \texttt{\tiny{[P1$\langle$GM]}}
                \texttt{You have attempted this same task before and did not succeed. These are the plans you wrote after those attempts.} \\
\texttt{Trial \#1: New plan: I incorrectly identified the causal structure as a simple fork with a confounding path, failing to recognize that the prompt explicitly defines a specific chain: Captain \$\textbackslash{}to\$ Private \$\textbackslash{}to\$ Prisoner and Captain \$\textbackslash{}to\$ Corporal \$\textbackslash{}to\$ Prisoner. In this structure, the Private is a **mediator** (on the causal pathway), not a confounder. Conditioning on a mediator (Method 2) blocks the indirect effect, preventing the estimation of the *total* effect of the Captain on the Prisoner. Method 1 (marginal correlation) correctly captures the total effect because there are no unmeasured confounders between the treatment and outcome in this self{-}contained world. Therefore, the correct answer is "Yes," as Method 1 is superior for assessing the total causal impact, while Method 2 suffers from overadjustment bias by blocking the very pathway we aim to measure.} \\
\\ 
\\ 
\texttt{You are an expert in causal inference. The following question is not a typical commonsense query, but rather a meticulously designed question created by a professor specializing in causal inference, intended to assess the students' mastery of the course content.} \\
\texttt{Imagine a self{-}contained, hypothetical world with only the following conditions, and without any unmentioned factors or causal relationships: The captain has a direct effect on the private and the corporal. The corporal has a direct effect on prisoner. The private has a direct effect on prisoner. Method 1: We look directly at how the captain correlates with prisoner in general. Method 2: We look at this correlation case by case according to the private. To understand how the captain affects prisoner, is it more correct to use the Method 1 than Method 2?} \\
\texttt{Start your answer with 'Yes' or 'No', followed by additional reasoning or evidence to support your explanation.} \\
            }
        }
    }
     \\ \\

    \theutterance \stepcounter{utterance}  
    & \multicolumn{4}{p{0.6\linewidth}}{
        \cellcolor[rgb]{0.561,0.737,0.561}{
            \makecell[{{p{\linewidth}}}]{
                \texttt{\tiny{[P1$\rangle$GM]}}
                \texttt{Yes} \\
\\ 
\texttt{In this self{-}contained hypothetical world, the causal structure is a directed acyclic graph (DAG) where the Captain (\$C\$) influences the Prisoner (\$P\$) through two distinct pathways:} \\
\texttt{1.  \$C \textbackslash{}to \textbackslash{}text\{Private\} \textbackslash{}to P\$} \\
\texttt{2.  \$C \textbackslash{}to \textbackslash{}text\{Corporal\} \textbackslash{}to P\$} \\
\\ 
\texttt{The Private acts as a **mediator** on the first pathway, transmitting the Captain's influence to the Prisoner. The Corporal acts as a mediator on the second pathway. There are no unmeasured confounders, and no back{-}door paths exist that would require adjustment to block bias.} \\
\\ 
\texttt{**Method 1** calculates the marginal correlation between the Captain and the Prisoner. In a causal framework without confounding, this marginal association captures the **total causal effect** of the Captain on the Prisoner. It includes both the direct effects (if any existed, though here the effect is fully mediated) and the indirect effects flowing through the Private and the Corporal.} \\
\\ 
\texttt{**Method 2** conditions on the Private (stratifying by the Private's status). By conditioning on the Private, we block the path \$C \textbackslash{}to \textbackslash{}text\{Private\} \textbackslash{}to P\$. This isolates only the effect of the Captain on the Prisoner that does *not* go through the Private. However, since the prompt defines the Captain's influence on the Prisoner as flowing *through* the Private (and the Corporal), conditioning on the Private removes a significant portion of the causal mechanism. This results in **overadjustment bias** (specifically, blocking a mediator), leading to an estimate of a partial effect rather than the total effect.} \\
\\ 
\texttt{Therefore, to understand "how the captain affects prisoner" in the sense of the total impact within this closed system, Method 1 is superior because it preserves the full causal chain, whereas Method 2 artificially severs the link via the Private, providing an incomplete and misleading picture of the Captain's total influence.} \\
            }
        }
    }
    & & \\ \\

    \theutterance \stepcounter{utterance}  
    & & & \multicolumn{2}{p{0.3\linewidth}}{
        \cellcolor[rgb]{0.9,0.9,0.9}{
            \makecell[{{p{\linewidth}}}]{
                \texttt{\tiny{[GM$|$GM]}}
                \texttt{Yes} \\
            }
        }
    }
    & & \\ \\

    \theutterance \stepcounter{utterance}  
    & & & \multicolumn{2}{p{0.3\linewidth}}{
        \cellcolor[rgb]{0.9,0.9,0.9}{
            \makecell[{{p{\linewidth}}}]{
                \texttt{\tiny{[GM$|$GM]}}
                \texttt{yes} \\
            }
        }
    }
    & & \\ \\

    \theutterance \stepcounter{utterance}  
    & & & \multicolumn{2}{p{0.3\linewidth}}{
        \cellcolor[rgb]{0.9,0.9,0.9}{
            \makecell[{{p{\linewidth}}}]{
                \texttt{\tiny{[GM$|$GM]}}
                \texttt{game\_result = WIN} \\
            }
        }
    }
    & & \\ \\

\end{supertabular}
}

\end{document}
